%% aidisclose.tex --- Generative AI disclosure checklist and statements
%%
%% Copyright (C) 2025 by João M. Lourenço
%%
%% This file may be distributed and/or modified under the conditions of
%% the LaTeX Project Public License, either version 1.3c of this license
%% or (at your option) any later version. The latest version of this
%% license is in:
%%
%%    http://www.latex-project.org/lppl.txt
%%
%% and version 1.3c or later is part of all distributions of LaTeX
%% version 2006/05/20 or later.
%%

%—————————————————————————————————————————————————————————————————
% AI Disclosure Customization
%.................................................................
%
% Preenchido em 2026-09-16 a partir da lista assinalada pelo autor.
% Três itens foram ACRESCENTADOS pelo assistente e estão marcados
% [ACRESCENTADO] com a razão: confirmar ou remover.
%—————————————————————————————————————————————————————————————————

% 1 — Conceptualization
\AIDactivate{c:idea}          % Geração de ideias
% \AIDactivate{c:obj}         % Definição do objetivo de investigação
\AIDactivate{c:rq}            % Formulação de questões de investigação e hipóteses
% \AIDactivate{c:feas}        % Avaliação de viabilidade e análise de risco
\AIDactivate{c:pre}           % Verificações preliminares (exploratórias) de hipóteses
\AIDactivate{c:sim}           % Debate simulado e teste de argumentos

% 2 — Literature Review
% \AIDactivate{l:srch}        % Pesquisa e Descoberta
\AIDactivate{l:sum}           % Resumo de fontes externas
\AIDactivate{l:map}           % Mapeamento de Conceitos e Sistematização
% \AIDactivate{l:pat}         % Análise de panorama de mercado/patentes
\AIDactivate{l:gaps}          % Identificação de Lacunas e Avaliação de Novidade
% \AIDactivate{l:trans}       % Compreensão de literatura multilingue

% 3 — Methodology
\AIDactivate{m:des}           % Desenho de investigação
\AIDactivate{m:proto}         % Desenvolvimento de protocolos experimentais ou de investigação
% \AIDactivate{m:meth}        % Seleção de métodos de investigação

% 4 — Software Development and Automation
\AIDactivate{s:gen}           % Geração de código
\AIDactivate{s:opt}           % [ACRESCENTADO] Refatoração e Otimização — o verificador de
                              % ancoragem, o runner e a reconciliação de quarentena foram
                              % reescritos várias vezes, não apenas gerados
\AIDactivate{s:debug}         % Depuração e Reparação
\AIDactivate{s:auto}          % Automação de processos
\AIDactivate{s:algs}          % [ACRESCENTADO] Design de algoritmos — detector de numeração de
                              % linhas, router de interleaving, correcção de localizador,
                              % orçamento de tentativas: são algoritmos desenhados, não só código
\AIDactivate{s:doc}           % Documentação de código e geração de comentários

% 5 — Data Management
% \AIDactivate{d:coll}        % Recolha de dados
\AIDactivate{d:val}           % Validação
\AIDactivate{d:cln}           % Limpeza de dados
\AIDactivate{d:cur}           % Curadoria e organização de dados
\AIDactivate{d:anl}           % Análise de dados
\AIDactivate{d:viz}           % Geração de figuras a partir de dados
\AIDactivate{d:rep}           % Verificações de reprodutibilidade e re-execução
\AIDactivate{d:lbl}           % [ACRESCENTADO] Assistência na rotulagem e anotação de dados —
                              % a extração do Stage C é exactamente isto: atribuição de
                              % categorias registadas a 41 artigos. É o que o item 10 da
                              % checklist PRISMA-ScR declara; omiti-lo aqui contradiria a tese
% \AIDactivate{d:syn}         % Geração de dados sintéticos
% \AIDactivate{d:anon}        % Apoio à desidentificação e anonimização
% \AIDactivate{d:trans}       % Transcrição de áudio para texto e diarização

% 6 — Visuals and Multimedia
\AIDactivate{v:gen}           % Geração de imagem/diagrama (sintética)
\AIDactivate{v:edit}          % Melhoria e Edição de Imagem
% \AIDactivate{v:chart}       % Design e layout de figuras

% 7 — Writing and Editing
\AIDactivate{w:draft}         % Redação de Texto
\AIDactivate{w:poly}          % Polimento e Edição
\AIDactivate{w:sum}           % Resumo de texto
% \AIDactivate{w:con}         % Formulação de conclusões
% \AIDactivate{w:tone}        % Adaptação e ajuste de tom emocional
\AIDactivate{w:tra}           % Tradução
\AIDactivate{w:ref}           % Formatação e Codificação de Referências
% \AIDactivate{w:prs}         % Comunicados de imprensa e materiais de divulgação
% \AIDactivate{w:title}       % Geração de Título

% 8 — Ethics Review
% \AIDactivate{e:bias}        % Análise de viés e avaliação de discriminação
% \AIDactivate{e:risk}        % Análise de risco ético
% \AIDactivate{e:comp}        % Monitorização de conformidade com padrões éticos
\AIDactivate{e:conf}          % Monitorização de confidencialidade de dados

% 9 — Quality Assurance
\AIDactivate{sup:qa}          % Revisão por Pares Simulada
% \AIDactivate{sup:trd}       % Identificação de tendências
\AIDactivate{sup:lim}         % Identificação de limitações
\AIDactivate{sup:rec}         % Sugestões Estratégicas
% \AIDactivate{sup:pub}       % Apoio à publicação


% --- GAI tools used ---
% CONFIRMAR: acrescentar ou remover conforme o que foi efectivamente usado.
\AIDtoolsUsed{Claude (Anthropic), Claude Code (Anthropic), ChatGPT (OpenAI)}

%—————————————————————————————————————————————————————————————————
% End-of AI Disclosure Customization
%—————————————————————————————————————————————————————————————————
